\scanpage{5}
le sous groupe
\[
\Gamma\subset\overline Q^*
\]
formé des $\lambda$ tels que l'on ait
\[
\left\{\begin{array}{l}
\lambda\unsure{\in}\ \text{entiers algébriques}\\
\lambda\ \text{une unité }\ell\text{-adique pour }\ell\ne p\\
|\lambda_i|=p^{\nu/2}\ \text{pour }\deleted{\unsure{\exists}}\ \nu\in\mathbf Z^+\ \text{convenable,}\\
\qquad\text{pour tt conjugué de }\lambda_i
\end{array}\right.
\]
Le groupe $G^0$ est défini par
\[
\Gamma/\Gamma_{\mathrm{tors}}=\Gamma/(\text{racines de l'unité})
\]
L'hom $G\mathbin{\unsure{-|\longrightarrow}^{\unsure{\epsilon}}}G_m$ correspond à l'hom
\diagram{\begin{tikzcd}
\mathbf Z \arrow[r] & \Gamma\\
1\arrow[r] & p
\end{tikzcd}}
et $\operatorname{Ker}(\epsilon|G^0)=\operatorname{Ker}(G^0-G_m)$ est défini par le groupe de caractères
\[
\Gamma/\{\text{racines de l'unité},p^\nu\}=\Gamma'
\]
qui est un groupe \deleted{\unsure{ent.}} dont le sous groupe de torsion est $\simeq\mathbf Z/2\mathbf Z$
\diagram{\begin{tikzcd}
\unsure{0}\arrow[r] & \mathbf Z/2\mathbf Z\arrow[r] & \Gamma'\arrow[r] & \Gamma'/\operatorname{tors}(\Gamma')\arrow[r] & 0\\
&1\arrow[r,mapsto]&\begin{array}{c}\text{classe de}\\p^{1/2}\text{ dans }\Gamma'\end{array}&&
\end{tikzcd}}
On a donc
\diagram{\begin{tikzcd}
\unsure{0}\arrow[r]&H^0\arrow[r]&H\arrow[r]&\mu_2\arrow[r]&0
\end{tikzcd}}
et $H^0$ est un \uline{\unsure{tor. compact}}. En fait, l'hom $i:G_m-G$ définit $i|\mu_2:\mu_2\longrightarrow H^0$, et donne un \uline{splitting}
\[
H\simeq H^0\times\mu_2
\]
\mspage{4}
